% SOURCE_STATEMENT: "For the domain-wall minimization problem, is the constrained infimum attained for every $0<\\theta\\leq\\pi/2$ and every $q\\in\\mathbb{Z}$, when $q$ is the degree or winding number?"
% Preparation and final self-audit: CONTRIBUTING.md prompts C and D; source/audit artifacts retained in the local physics_statements_2026_10_11 batch.
% !TeX program = lualatex
\documentclass[11pt]{article}
\usepackage[margin=25mm]{geometry}
\usepackage{fontspec}
\setmainfont{Latin Modern Roman}
\usepackage{amsmath,amssymb,mathtools}
\usepackage{unicode-math}
\setmathfont{Latin Modern Math}
\usepackage{xurl,hyperref}
\hypersetup{colorlinks=true,urlcolor=blue}
\setcounter{secnumdepth}{0}
\setlength{\parindent}{0pt}
\setlength{\parskip}{5pt}
\setlength{\emergencystretch}{3em}
\newcommand{\sref}[2]{\hyperlink{source:#1:#2}{[S#2]}}
\newcommand{\cataloguescope}[1]{\paragraph{Recorded target.}#1}
\newcommand{\R}{\mathbb R}
\newcommand{\C}{\mathbb C}
\newcommand{\N}{\mathbb N}
\newcommand{\Z}{\mathbb Z}
\newcommand{\Q}{\mathbb Q}
\newcommand{\F}{\mathbb F}
\newcommand{\eps}{\varepsilon}
\begin{document}
\section{Attainment of Constrained Domain-Wall Energies}\label{statement}

\subsection{Definitions and mathematical statement}

\paragraph{Shared notation.}$\mathbb N=\{1,2,\ldots\}$; the standard number fields have their usual meanings. All additional conventions are stated locally.


\paragraph{Definitions and conventions.}
In the report, $\Omega=\mathbb R\times(-1,1)$ has coordinates $(x_1,x_3)$, and magnetizations $m=(m_1,m_2,m_3)$ take values in $S^2$ with finite Dirichlet energy. Put $m'=(m_1,m_3)$ and $m_\theta^\pm=(\cos\theta,\pm\sin\theta,0)$. The class $\mathcal M_\theta$ has $\operatorname{div}m'=0$ in $\Omega$, $m_3=0$ on its boundary, and the source far-field condition $m(\pm\infty,\cdot)=m_\theta^\pm$. Its energy is $\int_\Omega|\nabla m|^2\,dx$. The two topological constraints are $D(m)=(2\pi)^{-1}\int_\Omega m\cdot(\partial_1m\times\partial_3m)\,dx$ and $W(m)=\pi^{-1}\int_\Omega\nabla m_1\times\nabla m_2\,dx$, where the last cross product is the determinant in the $(x_1,x_3)$ plane. The infimum is taken over $m\in\mathcal M_\theta$ with either $D(m)=q$ or $W(m)=q$.

\paragraph{Target retained from the canonical source catalogue.}
\paragraph{Catalogue wording retained for traceability.}
For the domain-wall minimization problem, is the constrained infimum attained for every $0<\theta\leq\pi/2$ and every $q\in\mathbb{Z}$, when $q$ is the degree or winding number?

\paragraph{Source scope and qualification.}
The energy, strip and both constraints were recovered from PDF page 21. The source first singles out $q\in\{-1,0,1\}$ and also asks the general integer case, so both are retained. Its negative-end square-integrability line displays a questionable plus sign inconsistent with its stated far-field value; that sign is not silently used to define an empty class. The asymptotic notation above retains the intended boundary-state statement, and a theorem must explicitly fix its convergence convention.
\subsection{Short English statement}
Does every allowed angle and integer winding class have an energy-minimizing domain wall?
\subsection{Sources}
\begin{itemize}
\item[\hypertarget{source:statement:1}{S1}] ULAM.AI and the UnsolvedMath Contributors, UnsolvedMath, catalogue entry. Catalogue curation under CC BY 4.0; definitions and English exposition added locally.
\url{https://huggingface.co/datasets/ulamai/UnsolvedMath}.
\item[\hypertarget{source:statement:2}{S2}] Mechanics of Materials: Mechanics of Interfaces and Evolving Microstructure (2016), Oberwolfach Reports, DOI 10.4171/owr/2016/17, pp. 797--867. Primary report accessed 2026-10-11; relevant PDF pages 21, 22.
\url{https://doi.org/10.4171/owr/2016/17}.
\item[\hypertarget{source:statement:3}{S3}] Döring--Ignat--Otto (catalogue-listed context; not independently consulted in this preparation).
\url{https://doi.org/10.4171/JEMS/464}.
\end{itemize}
\label{end:statement}
\end{document}
